\documentclass[11pt]{article}
\usepackage[letterpaper,margin=1in]{geometry}
\usepackage{amsmath}
\usepackage{amssymb}
\usepackage{booktabs}
\usepackage{tabularx}
\usepackage{longtable}
\usepackage[hidelinks]{hyperref}

\newcommand{\Rcal}{\mathcal{R}}
\newcommand{\Scal}{\mathcal{S}}
\newcommand{\Lcal}{\mathcal{L}}
\newcommand{\score}{\operatorname{score}}
\newcommand{\save}{\operatorname{save}}
\newcommand{\lp}{\operatorname{lp}}

\title{Revised Experiments for PatternBoost Geometric Extremal Search}
\author{Experiment design note}
\date{\today}

\begin{document}
\maketitle

\begin{abstract}
This note replaces the earlier broad ablation schedule for the PatternBoost
geometric extremal-search project.  The main observation from comparing the
new multi-level experiments with the earlier code is that better bounds came
from certificate-shaped search spaces, not from running a generic coordinate
mutation pipeline longer.  The revised experiment plan therefore uses
problem-specific representations and verifiers for the three current targets:
MISR LP gaps, unit-square line-stabbing gaps, and recursive guillotine
hardness.  Epsilon-net and representability-separation experiments are outside
the present scope.
\end{abstract}

\section{Revision principle}

The previous design treated the experiment as a large uniform matrix over
representations, local moves, and surrogate scores.  That matrix was useful for
debugging the harness, but it is not the right primary experiment for obtaining
meaningful mathematical bounds.  The comparison with the earlier code shows the
following pattern.

\begin{enumerate}
\item Unit-square progress came from a one-dimensional staircase certificate
that was lifted to a two-dimensional square-stabbing instance.
\item MISR progress came from changing the representation: endpoint-order
sequence pairs found stronger small LP gaps than direct coordinate mutation.
\item Guillotine progress should be formulated as a threshold obstruction
problem over large saved subsets, with nonseparable cores and witness-breaking
moves, rather than as a flat search over rectangle coordinates.
\end{enumerate}

Thus the new experimental question is not ``Which cell in a generic grid wins?''
The new question is:
\[
        \text{Which certificate family, verifier, and search operator
        produces the strongest auditable bound?}
\]

The revised experiments should report record-quality certificates and
diagnostic ablations of the specific mathematical mechanisms that produced
them.  Averages over many nominally identical random launches are not needed
for the manuscript.

\section{Global requirements}

Every reported result must satisfy the following requirements.

\begin{enumerate}
\item The reported value is recomputed by an exact verifier independent of the
training loop.
\item The certificate is standalone: it contains the geometry or combinatorial
template, objective components, solver status, and a canonical hash.
\item The paper distinguishes between search score and proof score.  Surrogate
values, sampled checks, and neural selection scores are never presented as
certified bounds.
\item Each problem has a known-baseline reproduction step before any new
search is reported.
\item The generic multi-level runs may be used as diagnostics, but not as the
main evidence for record bounds.
\end{enumerate}

\section{Experiment schedule}

\begin{center}
\begin{tabularx}{\textwidth}{p{0.18\textwidth}p{0.25\textwidth}X}
\toprule
Phase & Output & Purpose \\
\midrule
Artifact audit & Verified reproduction of known certificates & Confirm that the old strong constructions and the new verifier agree. \\
Targeted search & Candidate families and near-miss logs & Explore the certificate-shaped search spaces described below. \\
Mechanism ablation & Small tables comparing components inside one problem & Test which mathematical ingredient matters: representation, verifier, local move, or model proposal. \\
Record run & Best certificate per problem, plus verifier output & Produce final manuscript claims and certificate gallery. \\
Independent audit & Recomputed objective values from frozen certificates & Guard against logging, solver, or search-pipeline mistakes. \\
\bottomrule
\end{tabularx}
\end{center}

The schedule is deliberately organized by mathematical target rather than by a
uniform Cartesian product of implementation choices.  For each problem, the
primary run should be the one that searches the most structured family first.
Direct coordinate mutation is retained only as a sanity check.

\section{Problem 1: unit-square line stabbing}

\subsection{Objective}

For a finite family $\Scal$ of axis-parallel unit squares, the exact score is
\[
        \score_{\rm stab}(\Scal)=\frac{\tau(\Scal)}{\tau^*(\Scal)}.
\]
Here $\tau(\Scal)$ is the minimum number of horizontal and vertical lines
stabbing all squares, and $\tau^*(\Scal)$ is the LP optimum of the corresponding
finite set-cover relaxation over critical horizontal and vertical line
positions.

\subsection{Primary experiment: staircase interval certificates}

The primary unit-square experiment should be the staircase certificate search
from the older successful code.  A candidate is a one-dimensional interval
pattern with levels.  A rational measure on the line gives lower bounds
$a_r$ on the mass of intervals in level $r$.  The certificate proves that
fewer than a target number of lines cannot stab the two-dimensional product
instance whenever
\[
        a_r+a_s\geq 1
        \qquad\text{for all } r+s<k.
\]
The product instance includes $I\times J$ whenever the mass of $I$ plus the
mass of $J$ is at least one.  After scaling, these products become unit
squares.

This experiment should first reproduce the known certificates:
\[
        \frac{14}{9}=1.555\ldots,
        \qquad
        \frac{20}{13}=1.538\ldots.
\]
The $14/9$ certificate is the current unit-square baseline for the manuscript.

\subsection{Search variants}

The staircase search should vary the following ingredients.

\begin{enumerate}
\item Level count and target number of required lines.
\item Interval starts and common interval length.
\item Rational support points for the one-axis measure.
\item Deterministic scans over shifted-disjoint and near-staircase families.
\item PatternBoost proposals trained on high-scoring interval patterns, not on
raw square coordinates.
\end{enumerate}

The raw square-coordinate Axplorer environment remains a comparison baseline.
It is useful for confirming that a less structured search can find small
examples such as the $3/2$ pattern, but it should not be the primary search for
new record certificates.

\subsection{Verifier and certificate}

The unit-square certificate must contain:
\begin{enumerate}
\item the interval levels and common length;
\item the rational measure and support;
\item level piercing numbers;
\item all inequalities $a_r+a_s\geq 1$ used in the proof;
\item the generated product squares;
\item the compressed critical line universe;
\item exact LP and IP values for the finite square instance.
\end{enumerate}

The verifier should recompute both the symbolic staircase proof and the finite
LP/IP confirmation.  The manuscript table should report the exact ratio, the
number of intervals, number of product squares, LP value, IP value, and
certificate hash.

\section{Problem 2: MISR LP gap}

\subsection{Objective}

For a finite family $\Rcal$ of axis-parallel rectangles, the exact score is
\[
        \score_{\rm MISR}(\Rcal)=\frac{\alpha^*(\Rcal)}{\alpha(\Rcal)}.
\]
Here $\alpha(\Rcal)$ is the maximum number of pairwise-disjoint rectangles, and
$\alpha^*(\Rcal)$ is the standard clique LP optimum for the rectangle
intersection graph.

\subsection{Current baseline: endpoint-order sequence pairs}

The active MISR baseline is no longer the direct-coordinate search.  A candidate
is encoded by two double-occurrence sequences $H$ and $V$, each containing every
rectangle label exactly twice.  Rectangle $i$ receives its $x$-span from the two
positions of $i$ in $H$ and its $y$-span from the two positions of $i$ in $V$.
The decoded rectangle family is then sent to the ordinary exact MISR verifier.

This representation is the main useful difference from the earlier successful
code.  It searches endpoint orders rather than raw coordinates, so swaps,
block moves, reversals, label-pair swaps, motif refreshes, and depth-weighted
lifting all remain in a compact geometric search space.  The role of this
baseline is to test whether the endpoint-order inductive bias consistently
beats direct coordinate mutation.

\subsection{Primary search: sequence-pair PatternBoost}

The main MISR search should train and mutate in the endpoint-sequence space,
while keeping the exact verifier outside the training loop as the final
authority.  The preferred cell is:
\[
  \texttt{endpoint\_sequence\_pair}
  \times
  \texttt{sequence\_pair\_pivot}
  \times
  \texttt{exact\_lp\_gap\_pressure}.
\]
The exact surrogate is expensive, but for MISR it is scientifically preferable
to ranking candidates by a weak graph proxy when the target is a record LP gap.
Cheaper proxies remain in the matrix only as controls.

\subsection{Search tracks}

The MISR experiments should have three tracks.

\begin{enumerate}
\item \textbf{Sequence-pair track.}  Use double-occurrence endpoint sequences,
motif seeds, local endpoint-order moves, recombination, and lifting.  Report
only decoded rectangle certificates verified by the framework scorer.
\item \textbf{Direct-coordinate diagnostic.}  Keep direct coordinates as a
negative control.  It should explain the earlier $1.25$ plateau rather than
serve as the main search.
\item \textbf{Structured-program follow-up.}  If sequence-pair search stalls,
test triangle-free or bounded-program variants as additional structured
families.  These should be treated as follow-up, not as the first run.
\end{enumerate}

Four-port reproduction is explicitly outside the current active plan.  It can
be revived as a separate archival/certificate project, but the present MISR
experiments should focus on endpoint-sequence search and explicit finite
rectangle certificates.

\subsection{Verifier and certificate}

For an unweighted finite certificate, save:
\begin{enumerate}
\item rectangle coordinates;
\item intersection graph;
\item triangle-free status and witness if violated;
\item all maximal clique constraints, or all edge constraints in the
triangle-free case;
\item LP optimum $\alpha^*$ and primal solution;
\item exact maximum independent set and value $\alpha$;
\item solver status, runtime, and certificate hash.
\end{enumerate}

For sequence-pair candidates, also save the two sequences $H,V$, the decoded
rectangles, and the sequence mutation/lift history when available.  The
manuscript should distinguish the search representation from the certified
geometry: the proof object is still the decoded rectangle family.

\section{Problem 3: recursive guillotine hardness}

\subsection{Objective}

For a pairwise-disjoint family $\Rcal$ of axis-parallel rectangles, define
$\save(\Rcal)$ as the maximum number of rectangles that can be preserved by a
recursive guillotine cutting strategy.  The exact destroyed-fraction score is
\[
        \score_{\rm guill}(\Rcal)
        =
        1-\frac{\save(\Rcal)}{|\Rcal|}.
\]

The generic exact DP can compute this value for finalists.  However, it is not
the best search objective.  The stronger experimental formulation is a
threshold obstruction:
\[
        K=\left\lfloor\frac{|\Rcal|}{2}\right\rfloor+1.
\]
If there is no guillotine-separable subset of size $K$, then every guillotine
strategy destroys at least half the rectangles.

\subsection{Primary experiment: threshold obstruction search}

The guillotine search should use the threshold formulation as its main driver.
For each subset $S$, build the $x$-projection and $y$-projection overlap
graphs.  A vertical guillotine cut avoiding all rectangles in $S$ exists
exactly when the $x$-projection graph induced by $S$ is disconnected; the
horizontal case is analogous.  This gives a fast recursive oracle for testing
whether $S$ is guillotine-separable.

The search score is the fraction of tested $K$-subsets that are
nonseparable.  A certified pass has zero separable $K$-subsets.  Near-misses
with very few separable witnesses are valuable and should be retained for
local improvement.

\subsection{Local moves}

Use local moves that directly attack separable witnesses.
\begin{enumerate}
\item If a separable witness splits along a vertical gap, stretch or move
rectangles so that the corresponding $x$-projection components become linked.
\item If a witness splits along a horizontal gap, do the analogous move in the
$y$-projection.
\item Mine small nonseparable cores and preserve them during mutation.
\item Assemble recursive arrangements by placing hard cores so that every
large subset contains at least one nonseparable core.
\item Use long thin blockers, pinwheel motifs, brick-wall motifs, and
skew-grid templates as structured starting families.
\end{enumerate}

This is different from optimizing only the best first cut.  First-cut hardness
is a useful diagnostic, but the final claim must be about recursive
separability or exact saved count.

\subsection{Verifier and certificate}

The guillotine certificate should have two layers.

\begin{enumerate}
\item \textbf{Threshold certificate.}  Prove that no $K$-subset is
guillotine-separable, either by exact enumeration or by a nonseparable-core
cover certificate.
\item \textbf{Exact-score certificate.}  For smaller finalists, run the full
recursive DP and report $\save(\Rcal)$ and the destroyed fraction.
\end{enumerate}

For final reported arrangements, save the rectangles, $K$, separability oracle
status, number of checked subsets, any separable witness if the candidate
fails, nonseparable cores, exact DP value when computed, and certificate hash.

\section{Controls and comparisons}

The controls should compare mathematical mechanisms, not repeated launches.
Use the following comparisons.

\begin{enumerate}
\item \textbf{Structured versus direct representation.}  Compare the primary
certificate-shaped representation against direct coordinate mutation.
\item \textbf{Model proposal versus local search.}  Run the same verifier and
local moves with and without PatternBoost proposals.
\item \textbf{Surrogate relevance.}  Compare the search surrogate to exact
scores on audited candidates and report rank correlation or clear failure
cases.
\item \textbf{Known-baseline reproduction.}  Every problem should first
reproduce the strongest available old construction or explain why no such
construction exists.
\item \textbf{Near-miss analysis.}  For failed guillotine and MISR candidates,
report the obstruction that prevented certification: separable witness,
triangle, solver timeout, large independent set, or collapsed geometry.
\end{enumerate}

\section{Manuscript tables}

The main manuscript should contain compact tables of certified results, not
large random-launch matrices.

\begin{center}
\begin{tabularx}{\textwidth}{p{0.2\textwidth}p{0.24\textwidth}p{0.2\textwidth}X}
\toprule
Problem & Primary search family & Certified objective & Baseline to reproduce \\
\midrule
Unit-square stabbing & Staircase interval certificate and product lift & $\tau/\tau^*$ & $14/9$ certificate \\
MISR & Endpoint-order sequence pairs with exact LP-gap audit & $\alpha^*/\alpha$ & Endpoint-sequence smoke certificate \\
Guillotine & Threshold nonseparable-subset search & destroyed fraction or no separable $K$-subset & Current small obstruction motifs; no stronger old certificate found \\
\bottomrule
\end{tabularx}
\end{center}

For each final result, report:
\begin{enumerate}
\item exact ratio or threshold statement;
\item instance size;
\item objective components;
\item verifier status;
\item runtime of the verifier;
\item certificate hash;
\item path to rendering or gallery figure.
\end{enumerate}

The appendix may include diagnostic tables for direct-coordinate experiments,
but these should be framed as evidence that generic search is insufficient for
record bounds.

\section{Artifact layout}

The final artifact should organize outputs by problem and certificate family.
A recommended layout is:
\begin{verbatim}
artifacts/
  unit_square/
    staircase_14_9/
      certificate.json
      verify_certificate.py
      rendered_instance.pdf
    staircase_search/
      run_config.json
      candidates.jsonl
      audited_summary.csv
  misr/
    four_port/
      certificate.json
      verify_four_port.py
      verification_output.txt
    triangle_free_programs/
      programs.jsonl
      best_certificates/
      audited_summary.csv
  guillotine/
    threshold_search/
      best_instance.rect
      threshold_certificate.json
      separable_witnesses.jsonl
      nonseparable_cores.jsonl
    exact_dp_audit/
      dp_certificate.json
\end{verbatim}

Each verifier should run without training logs.  Training logs are useful for
reproducibility, but certificate verification must not depend on them.

\section{What not to run}

The following experiments should not be part of the immediate manuscript plan.

\begin{enumerate}
\item The full generic $3\times 3\times 3$ matrix as the primary evidence.
\item Long direct-coordinate MISR runs whose best result remains near small
cycle-like plateaus.
\item Guillotine runs scored only by first-cut hardness.
\item Epsilon-net and graph-representability experiments, unless the paper is
expanded beyond the three current geometric objectives.
\item Any result whose exact verifier times out or whose certificate cannot be
recomputed from a standalone file.
\end{enumerate}

\section{Summary}

The revised experimental strategy is certificate-first.  Unit-square stabbing
should use staircase interval certificates.  MISR should use endpoint-order
sequence pairs as the primary search space, with direct coordinates retained
only as a diagnostic baseline.  Guillotine hardness should use the threshold
nonseparable-subset formulation and core-cover certificates, with full DP only
for finalists.  Generic PatternBoost remains useful as a proposal mechanism
and diagnostic baseline, but the mathematical search space must encode the
proof pattern we want the system to discover.

\end{document}
